\chapter{Methodology and Proposed Approach}
\label{chap:methodology-proposed-approach}

This chapter explains how the work is organized to answer the problem. It should justify the approach before describing the detailed design.

\section{Engineering Approach}
Describe the development or research process used in the project: iterative development, experimental protocol, design science, agile workflow, modeling approach, or another suitable method.

\placeholderfigure{methodology-workflow.pdf}{Methodology workflow placeholder}{Insert the main workflow followed during the project.}

Avoid writing a calendar diary. Report the method, not every daily activity.

\section{Design Logic}
Explain the reasoning that connects the requirements to the proposed solution. Show why the chosen approach fits the context and constraints.

Avoid presenting design decisions as arbitrary. Each important choice should answer a need, constraint, risk, or lesson from the literature review.

\section{Tools and Technologies}
List the tools, languages, frameworks, datasets, devices, or platforms used. For each one, state its role in the project.

Avoid tool catalogs. Do not include logos or marketing claims; include only technologies that matter to implementation or validation.

\section{Justification of Choices}
Compare reasonable alternatives and explain why the selected option was retained. Mention trade-offs honestly.

\begin{table}[H]
    \centering
    \caption{Choice justification template}
    \label{tab:choice-justification-template}
    \small
    \begin{tabular}{L{0.24\textwidth}L{0.30\textwidth}L{0.30\textwidth}}
        \toprule
        \textbf{Decision} & \textbf{Reason retained} & \textbf{Trade-off or risk} \\
        \midrule
        Replace with a choice. & Explain the evidence or constraint. & State the cost, limitation, or mitigation. \\
        \bottomrule
    \end{tabular}
\end{table}

\section{Conclusion}
Summarize the adopted methodology and prepare the reader for the analysis and design chapter.